% !TeX program = xelatex
% !TeX options = -shell-escape
% !TEX encoding = UTF-8

% --------------------------------------------------------------
% This is all preamble stuff that you don't have to worry about.
% Head down to where it says "Start here"
% --------------------------------------------------------------

\documentclass[12pt]{article}
\usepackage[margin=2cm]{geometry} % 頁面邊距設置
\usepackage{amsmath,amssymb,amsthm} % 數學相關套件
\usepackage{graphicx} % 圖形插入
\usepackage{booktabs} % 表格線條
\usepackage[AutoFakeBold,AutoFakeSlant]{xeCJK} % 中文支援
\usepackage{fancyhdr} % 頁首頁尾設置
\usepackage{xcolor,changepage,mdframed} % 顏色與框線相關
\usepackage{setspace} % 行距
\usepackage{minted} % 程式碼高亮（需安裝 Python 套件 pygments）
\usepackage{caption}
\usepackage{tabularx}
\usepackage{longtable}
\usepackage{array}
\usepackage{ulem}
\usepackage[shortlabels]{enumitem} % 列表格式
\usepackage[colorlinks=true,linkcolor=black,urlcolor=blue,citecolor=blue]{hyperref}
\usepackage{multicol}
\usepackage{parskip}
\usepackage{paracol}
\setlength{\parskip}{6pt plus 2pt minus 1pt} % 段落間距
\graphicspath{{images}}
\usepackage{tikz}
\usetikzlibrary{arrows.meta, positioning}
% 字體設定 (xeCJK)
\setCJKmainfont{黑體-繁}
\setCJKmonofont{Noto Sans Mono CJK TC}

% 頁首/頁尾設定 (fancyhdr)
\pagestyle{fancy}
\fancyhead[L]{新北捷運公司} % 可以用 \leftmark 取代
\fancyhead[R]{求解器數學模型}
\fancyfoot[C]{\thepage}
\renewcommand{\headrulewidth}{0.4pt}
\renewcommand{\footrulewidth}{0pt}
\setlength{\headheight}{15pt}
\usepackage{siunitx}
\sisetup{
  round-mode=places,
  round-precision=6,
  table-number-alignment = center,
  detect-mode,
  detect-weight=true,
  detect-family=true
}

\captionsetup{width=0.45\textwidth}

% tex-fmt: off
% 自訂框線環境
\newmdenv[
  linewidth=2pt, linecolor=gray, backgroundcolor=gray!10,
  topline=false, bottomline=false, rightline=false,
  skipabove=10pt, skipbelow=10pt, innerleftmargin=10pt,
  innerrightmargin=10pt, innertopmargin=5pt, innerbottommargin=5pt
]{myquote}

\newenvironment{mycolorbox}[1][gray]{
  \renewmdenv[
    linewidth=1pt, linecolor=#1, backgroundcolor=#1!15,
    skipabove=10pt, skipbelow=10pt, innerleftmargin=10pt,
    innerrightmargin=10pt, innertopmargin=10pt, innerbottommargin=10pt
  ]{myquote}
  \begin{myquote}
}{\end{myquote}}

\newtheorem*{theorem}{Theorem}
\newenvironment{lemma}[2][Lemma]{\begin{trivlist}
\item[\hskip \labelsep {\bfseries #1}\hskip \labelsep {\bfseries #2.}]}{\end{trivlist}}
\newenvironment{exercise}[2][Exercise]{\begin{trivlist}
\item[\hskip \labelsep {\bfseries #1}\hskip \labelsep {\bfseries #2.}]}{\end{trivlist}}
\newenvironment{problem}[2][Problem]{\smallskip\begin{myquote}\begin{trivlist}
\item[\hskip \labelsep {\bfseries #1}\hskip \labelsep {\bfseries #2}]}{\end{trivlist}\end{myquote}}
\newenvironment{question}[2][Question]{\begin{trivlist}
\item[\hskip \labelsep {\bfseries #1}\hskip \labelsep {\bfseries #2.}]}{\end{trivlist}}
\newenvironment{corollary}[2][Corollary]{\begin{trivlist}
\item[\hskip \labelsep {\bfseries #1}\hskip \labelsep {\bfseries #2.}]}{\end{trivlist}}
\newenvironment{solution}{\begin{proof}[Solution]}{\end{proof}}
\newenvironment{answer}{\begin{proof}[Ans]}{\end{proof}}

% 樣式設定

\setminted{frame=lines, linenos, framesep=2mm, breaklines=true, breakanywhere=true, autogobble=true}

\linespread{1.1}
% \usemintedstyle{xcode}
\AtBeginEnvironment{minted}{\vspace{-10pt}} % 可根據需求調整間距數值
\setminted{breaklines=true, breakanywhere=true, autogobble=true}
\newcolumntype{Y}{>{\centering\arraybackslash}X}
\renewcommand\tabularxcolumn[1]{m{#1}} % 垂直置中
\setlength{\extrarowheight}{2pt}       % 多一點行距
\usepackage{ragged2e}
\renewcommand{\arraystretch}{1.2}
\renewcommand\tabularxcolumn[1]{>{\RaggedRight\arraybackslash}p{#1}}
\newcolumntype{P}[1]{>{\centering\arraybackslash}p{#1}}

% tex-fmt: on
\begin{document}
\thispagestyle{fancy}

%------------------------------------------------------------
%                         Start here
% -----------------------------------------------------------

\begin{center}
  {\large \bfseries 新北捷運公司人員自動排班系統} \\
  \vspace{1em}
  {\Large \bfseries 求解器數學模型} \\
\end{center}

\vspace{1em}

\begin{tabular}{@{}l l}
  \textbf{Intern:} & 吳亞倫 \\[-5pt]
  \textbf{ID:} & 1508S0342  \\[-5pt]
  \textbf{Department:} & 企劃處
\end{tabular}

\bigskip
\section{Problem Statement}

本研究建立站務人員 M 與檢修人員 T 的月排班數學模型。給定目標月份、
延伸日期、人員資料、歷史班表、指定休假、固定勤務、八週休假週期及各單位
營運需求，模型決定每位人員每日的正常班別、一般休假 \(R\)、國定假日休假
\(R1\) 或固定勤務 \(X\)。

模型先滿足所有硬性限制，再依群組順序改善指定休假、月休分布、工作區段、
班別安排及公平性。目標月份之後固定建立 \(K\) 日延伸日期，用來檢查月底安排
是否能合法延續；延伸日期不計入目標月份統計與軟性目標。

\section{Methodology}

\subsection{Optimization Structure}

令 \(\mathcal{F}\) 表示全部硬性限制所形成的可行域。每一條軟性規則定義一個
非負違反量
\[
  \Phi_j(\mathbf{x})\in\mathbb{Z}_{\geq0}.
\]
同一優先群組內的規則以加權和組成
\[
  J_p(\mathbf{x})
  =
  \sum_{j\in\mathcal{J}_p}\omega_j\Phi_j(\mathbf{x}),
\]
不同群組之間採字典序最佳化：
\[
  \operatorname{lexmin}_{\mathbf{x}\in\mathcal{F}}
  \left(J_1(\mathbf{x}),J_2(\mathbf{x}),\ldots,J_P(\mathbf{x})\right).
\]
亦即先最小化 \(J_1\)，固定其最優值後再最小化 \(J_2\)，依此類推。群組權重
只表示同一群組內不同違反量的交換比例。

為便於數學模型與規格書交叉查閱，規則使用短 ID：\texttt{CH}/\texttt{CS}
表示 M、T 共通的硬性／軟性規則，\texttt{MH}/\texttt{MS} 表示 M 規則，
\texttt{TH}/\texttt{TS} 表示 T 規則。ID 僅存在於文件並識別規則本身，不表示
優先順序；程式以有意義的英文函數名稱對應公式，不保存 Rule ID。

\subsection{The M Problem}

\subsubsection{Input and Output}

M 模型的輸入包括：
\begin{itemize}
  \item 目標月份與固定延伸天數 \(K\)；
  \item M 人員及其所屬車站；
  \item 指定休假 \(R^*\)、固定勤務 \(X\) 與其他固定指派；
  \item 八週週期的一般 \(R\) 與 \(R1\) 額度；
  \item 目標月份的一般 \(R\) 與 \(R1\) 基準數量；
  \item 月初前連續七個日曆日的歷史班表；
  \item 建模開始前各人員、各八週週期已使用的 R 與 R1 數量；
  \item 國定假日集合。星期六與星期日由日期本身決定。
\end{itemize}

模型輸出每位 M 人員在目標月份內的實際狀態。正常班需同時指出車站與班別；
滿足 \(R^*\) 時，輸出仍保留申請標記，並同時記錄其底層實際使用的是 \(R\)
或 \(R1\)。延伸日期只屬於模型內部決策，不列入正式輸出。

\subsubsection{Sets and Indices}

\begin{description}
  \item[\(\mathcal{E}^{M}\)] M 人員集合，索引為 \(e\)。
  \item[\(\mathcal{D}^{\mathrm{tar}}\)] 目標月份內的日期集合，索引為 \(d\)。
  \item[\(\mathcal{D}^{+}\)] 目標月份之後固定 \(K\) 日的延伸日期集合。
  \item[\(\mathcal{D}\)] 完整建模日期集合，
    \(\mathcal{D}=\mathcal{D}^{\mathrm{tar}}\cup\mathcal{D}^{+}\)。
  \item[\(\mathcal{L}\)] 車站集合，索引為 \(l\)。
  \item[\(\mathcal{S}\)] 正常班別集合，
    \(\mathcal{S}=\{\mathrm{E},\mathrm{A},\mathrm{N}\}\)，分別代表早、午、夜班。
  \item[\(\mathcal{C}\)] 八週休假週期集合，索引為 \(c\)。
  \item[\(\mathcal{G}\)] 車站群組集合。
  \item[\(\mathcal{D}^{wd}\)] 平日集合，即星期一至星期五且不是國定假日的日期。
  \item[\(\mathcal{D}^{hol}\)] 假日集合，即星期六、星期日或國定假日的日期。
  \item[\(\mathcal{L}^{ext}\)] 允許使用外派人力的車站集合。
\end{description}

車站群組固定為
\[
  \begin{aligned}
    \mathcal{L}_{A}&=\{\mathrm{LB01},\mathrm{LB02},\mathrm{LB03}\},\\
    \mathcal{L}_{B}&=\{\mathrm{LB04},\mathrm{LB05},\mathrm{LB06}\},\\
    \mathcal{L}_{C}&=\{\mathrm{LB07},\mathrm{LB08},\mathrm{LB09}\},\\
    \mathcal{L}_{D}&=\{\mathrm{LB10},\mathrm{LB11},\mathrm{LB12}\}.
  \end{aligned}
\]
允許外派的車站為
\[
  \mathcal{L}^{ext}=\{\mathrm{LB02},\mathrm{LB04},\mathrm{LB11}\}.
\]

\subsubsection{Parameters}

對每位人員 \(e\)，令 \(h_e\in\mathcal{L}\) 為所屬車站，\(g(l)\) 為車站
\(l\) 的群組，並定義合法工作車站
\[
  \mathcal{A}_e
  =
  \{l\in\mathcal{L}:g(l)=g(h_e)\}.
\]

令 \(b_{dls}\in\mathbb{Z}_{\geq0}\) 為日期 \(d\)、車站 \(l\)、班別 \(s\)
所需人數。M 的固定需求為所有車站每日早班與午班各一人，且
\(\mathrm{LB01}\)、\(\mathrm{LB06}\)、\(\mathrm{LB08}\)、\(\mathrm{LB12}\)
每日夜班各一人；平日、週末與國定假日使用相同需求。

M 的班別時間為：早班 \(06{:}30\)--\(14{:}30\)、午班
\(14{:}20\)--\(22{:}20\)、夜班 \(22{:}00\)--翌日 \(07{:}00\)。夜班歸屬
其開始日期。

令 \(p_{ed}\in\{0,1\}\) 表示人員 \(e\) 是否在日期 \(d\) 提出指定休假
\(R^*\)。\(R^*\) 不是獨立休假類型；若申請被滿足，該日仍須實際使用
\(R\) 或 \(R1\)。令 \(z_{ed}\in\{0,1\}\) 表示該日是否具有固定勤務 \(X\)，
且每筆 \(X\) 具有固定的實際開始與結束時間。

對每個八週週期 \(c\)，令 \(G_c\) 為一般 \(R\) 額度，令 \(H_c\) 為 \(R1\)
額度；目前一般額度通常為 \(G_c=16\)。令
\(\overline{R}_{ec}\)、\(\overline{R1}_{ec}\) 分別為建模區間開始前，週期
\(c\) 已使用的一般 \(R\) 與 \(R1\) 數量。令 \(n\)、\(m\) 分別為目標月份
的一般 \(R\) 與 \(R1\) 基準數量。

月初前連續七日的歷史班表用來重建法規連續日數、尚未結束的實際工作區段、
尚未結束的同班別區塊與最近一次有效正常班別。八週 R/R1 已使用數量由獨立
輸入提供，不從七日歷史推算；公平性不使用歷史累積值。

\subsubsection{Variables}

主要決策變數為
\[
  x^{M}_{edls}\in\{0,1\},
\]
其中 \(x^{M}_{edls}=1\) 表示人員 \(e\) 在日期 \(d\) 於車站 \(l\) 上班別
\(s\)。休假變數為
\[
  r^{M}_{ed},\ r^{M,1}_{ed}\in\{0,1\},
\]
分別表示一般 \(R\) 與 \(R1\)。外派班位變數為
\[
  a_{dls}\in\mathbb{Z}_{\geq0}.
\]

定義衍生變數
\[
  y^{M}_{eds}=\sum_{l\in\mathcal{A}_e}x^{M}_{edls},
  \qquad
  w^{M}_{ed}=\sum_{s\in\mathcal{S}}y^{M}_{eds}+z_{ed},
  \qquad
  \rho^{M}_{ed}=r^{M}_{ed}+r^{M,1}_{ed},
\]
其中 \(w^{M}_{ed}=1\) 表示實際工作日，\(\rho^{M}_{ed}=1\) 表示實際休假日。
另令
\[
  u_{ed}
  =
  \sum_{s\in\mathcal{S}}
  \sum_{\substack{l\in\mathcal{A}_e\\l\neq h_e}}x^{M}_{edls}
\]
表示該日是否跨站支援。

為避免混淆，下列狀態變數具有不同用途：

\begin{center}
\small
\begin{tabularx}{\textwidth}{p{2.6cm}X X}
  \toprule
  狀態 & 計數內容與中斷條件 & 用途 \\
  \midrule
  \(q^{law}_{ed}\) & 自最近一次一般 \(R\) 後經過的排班日數；只有一般 \(R\) 歸零，\(R1\) 仍增加 & \texttt{CH02} 法規性最多連續六日 \\
  \(c^{work}_{ed}\) & 連續實際工作日數；正常班與 \(X\) 增加，\(R\) 或 \(R1\) 歸零 & \texttt{CS04} 工作與休假區段長度 \\
  \((\lambda_{ed},\ell_{ed})\) & 最近有效正常班別及其連續出現次數；\(R\)、\(R1\)、\(X\) 都略過，遇到不同正常班才結束 & \texttt{MS03} 同班別區塊 \\
  \(\gamma_{ed}\) & 最近一次實際休假後出現的最近正常班別；\(R\)、\(R1\) 清除，\(X\) 略過 & \texttt{MS06} 換班前是否有休假 \\
  \(\eta_{ed}\) & 全序列中最近一次有效正常班別；\(R\)、\(R1\)、\(X\) 都略過 & \texttt{MS07} 班別輪轉方向 \\
  \bottomrule
\end{tabularx}
\normalsize
\end{center}

另定義三日模式指示變數
\(\pi^{NE}_{ed},\pi^{NA}_{ed}\in\{0,1\}\)，分別表示夜--休--早與夜--休--午。

\subsubsection{Constraints}

\paragraph{Daily unique assignment (\texttt{CH01})}
每位人員每天恰好具有一個實際狀態：
\[
  \sum_{l\in\mathcal{A}_e}\sum_{s\in\mathcal{S}}x^{M}_{edls}
  +r^{M}_{ed}+r^{M,1}_{ed}+z_{ed}=1,
  \qquad
  \forall e\in\mathcal{E}^{M},\ d\in\mathcal{D}.
\]
固定正常班、固定 \(R\)、固定 \(R1\) 與固定 \(X\) 直接固定對應變數。
\(R^*\) 只表示休假偏好，不固定底層休假類型。

\paragraph{Legal working stations (\texttt{MH01})}
對所有 \(l\notin\mathcal{A}_e\)，要求
\[
  x^{M}_{edls}=0.
\]
因此人員只能在本站或同群組車站工作。

\paragraph{Coverage (\texttt{MH02})}
每個必要班位必須完全補足：
\[
  \sum_{e\in\mathcal{E}^{M}}x^{M}_{edls}+a_{dls}=b_{dls},
  \qquad
  \forall d\in\mathcal{D},\ l\in\mathcal{L},\ s\in\mathcal{S}.
\]

\paragraph{External-station restriction (\texttt{MH03})}
對所有 \(l\notin\mathcal{L}^{ext}\)，要求
\[
  a_{dls}=0.
\]
外派只代表外部人力補足一個班位，不對應任何 M 人員列。

\paragraph{Minimum eleven-hour rest (\texttt{CH03})}
令 \(\mathcal{I}^{M}\) 為所有實際休息時間少於十一小時的正常班組合。對任一
\(((d,s),(d',s'))\in\mathcal{I}^{M}\)，要求
\[
  y^{M}_{eds}+y^{M}_{ed's'}\leq1,
  \qquad
  \forall e\in\mathcal{E}^{M}.
\]
若正常班與固定 \(X\) 的實際區間不足十一小時，該正常班變數固定為 0；兩筆固定
\(X\) 若彼此重疊或間隔不足十一小時，輸入本身不可行。此規則只比較實際工作
結束與下一次實際工作開始，與班別名稱是否相同無關。

\paragraph{At most six days without a general rest (\texttt{CH02})}
此規則是法規性連續日數限制。對每日 \(d\)，要求
\[
  r^{M}_{ed}=1
  \ \Longrightarrow\ 
  q^{law}_{ed}=0,
\]
以及
\[
  r^{M}_{ed}=0
  \ \Longrightarrow\ 
  q^{law}_{ed}=q^{law}_{e,d-1}+1.
\]
模型限制
\[
  0\leq q^{law}_{ed}\leq6.
\]
因此正常班、\(X\) 與 \(R1\) 都會使計數增加；只有一般 \(R\) 會歸零。
例如 \(\mathrm{E},\mathrm{E},\mathrm{E},R1,\mathrm{A},\mathrm{A},\mathrm{A}\)
會得到計數 \(1,2,3,4,5,6,7\)，故不合法。若 \(R^*\) 以 \(R\) 滿足則歸零，
以 \(R1\) 滿足則不歸零。

\paragraph{Eight-week rest quotas (\texttt{CH04})}
對每個週期 \(c\)，令
\[
  \mathcal{D}_c=c\cap\mathcal{D},
  \qquad
  F_c=\left|\{d\in c:d>\max\mathcal{D}\}\right|,
\]
其中 \(F_c\) 是模型區間結束後、週期內尚未建模的日期數。

若模型已涵蓋週期 \(c\) 的結束日，則
\[
  \overline{R}_{ec}
  +\sum_{d\in\mathcal{D}_c}r^{M}_{ed}
  =G_c,
  \qquad
  \overline{R1}_{ec}
  +\sum_{d\in\mathcal{D}_c}r^{M,1}_{ed}
  =H_c.
\]

若模型尚未涵蓋週期結束日，則兩種額度均不得超用：
\[
  \overline{R}_{ec}+\sum_{d\in\mathcal{D}_c}r^{M}_{ed}\leq G_c,
  \qquad
  \overline{R1}_{ec}+\sum_{d\in\mathcal{D}_c}r^{M,1}_{ed}\leq H_c,
\]
並且剩餘日期必須分別足以補足兩種額度：
\[
  \overline{R}_{ec}+\sum_{d\in\mathcal{D}_c}r^{M}_{ed}+F_c\geq G_c,
\]
\[
  \overline{R1}_{ec}+\sum_{d\in\mathcal{D}_c}r^{M,1}_{ed}+F_c\geq H_c.
\]
由於同一日不能同時使用 \(R\) 與 \(R1\)，另需滿足聯合可補足性：
\[
  \overline{R}_{ec}+\overline{R1}_{ec}
  +\sum_{d\in\mathcal{D}_c}\left(r^{M}_{ed}+r^{M,1}_{ed}\right)
  +F_c
  \geq G_c+H_c.
\]
\(R1\) 可安排在週期內任意日期，不必安排於國定假日當日。

\paragraph{Fixed assignments and historical boundary (\texttt{CH05})}
所有固定指派均不得由模型改寫。歷史班表決定建模區間第一日前的
\(q^{law}\)、\(c^{work}\)、\((\lambda,\ell)\)、\(\gamma\)、\(\eta\) 及週期
累積值，這些值作為第一日狀態轉移的固定前置狀態。

\paragraph{Requested rest (\texttt{CS01})}
只要指定休假日實際使用 \(R\) 或 \(R1\)，即視為申請被滿足。未滿足數為
\[
  \Phi^{M,R^*}
  =
  \sum_{\substack{e\in\mathcal{E}^{M},\ d\in\mathcal{D}^{\mathrm{tar}}\\p_{ed}=1}}
  \left(1-r^{M}_{ed}-r^{M,1}_{ed}\right).
\]

\paragraph{Monthly rest distribution (\texttt{CS02}, \texttt{CS03})}
令
\[
  Q^{M,R}_e=\sum_{d\in\mathcal{D}^{\mathrm{tar}}}r^{M}_{ed},
  \qquad
  Q^{M,R1}_e=\sum_{d\in\mathcal{D}^{\mathrm{tar}}}r^{M,1}_{ed}.
\]
定義
\[
  f_t(q)=\left[\max\{0,|q-t|-1\}\right]^2.
\]
則一般 \(R\) 與 \(R1\) 的月內分布違反量分別為
\[
  \Phi^{M,month,R}
  =
  \sum_{e\in\mathcal{E}^{M}}f_n\!\left(Q^{M,R}_e\right),
\]
\[
  \Phi^{M,month,R1}
  =
  \sum_{e\in\mathcal{E}^{M}}f_m\!\left(Q^{M,R1}_e\right).
\]
偏離基準 0 或 1 日不懲罰；偏離 2、3、4 日時，懲罰依序為 1、4、9。
此為月內分布偏好，不取代 \texttt{CH04} 的八週精確額度。

\paragraph{External staffing (\texttt{MS01})}
\[
  \Phi^{M,ext}
  =
  \sum_{d\in\mathcal{D}^{\mathrm{tar}}}
  \sum_{l\in\mathcal{L}^{ext}}
  \sum_{s\in\mathcal{S}}a_{dls}.
\]

\paragraph{Non-home-station assignments (\texttt{MS02})}
\[
  \Phi^{M,home}
  =
  \sum_{e\in\mathcal{E}^{M}}
  \sum_{d\in\mathcal{D}^{\mathrm{tar}}}u_{ed}.
\]

\paragraph{Continuous work-streak length (\texttt{CS04})}
此規則控制實際工作與休假的分段品質，不是 \texttt{CH02} 的法規計數。
實際工作區段是連續曆日皆為正常班或 \(X\) 的最大區段；一般 \(R\) 與
\(R1\) 都會結束區段。定義區段長度懲罰
\[
  D(L)
  =
  \begin{cases}
    4, & L=1,\\
    2, & L=2,\\
    1, & L=3,\\
    0, & 4\leq L\leq5,\\
    2(L-5), & L\geq6.
  \end{cases}
\]
令 \(\mathcal{B}^{work}_e\) 為人員 \(e\) 的最大實際工作區段集合，並承接月初
未結束區段。只對結束日在目標月份內的區段計分：
\[
  \Phi^{M,streak}
  =
  \sum_{e\in\mathcal{E}^{M}}
  \sum_{B\in\mathcal{B}^{work}_e:\ \max B\in\mathcal{D}^{\mathrm{tar}}}
  D(|B|).
\]

\paragraph{Same-shift block length (\texttt{MS03})}
此規則控制同一正常班別維持多久，用來減少生理時段頻繁切換。建立正常班別
序列時，\(R\)、\(R1\) 與 \(X\) 全部略過；只有下一個有效正常班別不同時，
原區塊才結束。令 \(\mathcal{B}^{shift}_e\) 為相同正常班別的最大區塊集合，
並承接月初未結束區塊，則
\[
  \Phi^{M,block}
  =
  \sum_{e\in\mathcal{E}^{M}}
  \sum_{B\in\mathcal{B}^{shift}_e:\ B\text{ 於目標月份內結束}}
  D(|B|).
\]
第一個正常班只建立第一個區塊，不產生前置區塊懲罰。

\paragraph{Night--rest--early pattern (\texttt{MS04})}
對所有完全位於目標月份內的三日視窗，令
\[
  \pi^{NE}_{ed}=1
  \Longleftrightarrow
  y^{M}_{ed\mathrm{N}}=1,
  \ \rho^{M}_{e,d+1}=1,
  \ y^{M}_{e,d+2,\mathrm{E}}=1.
\]
總違反量為
\[
  \Phi^{M,NE}=\sum_{e,d}\pi^{NE}_{ed}.
\]
此規則特別避免夜班後只休一日即接早班。

\paragraph{Night--rest--afternoon pattern (\texttt{MS05})}
同理，令
\[
  \pi^{NA}_{ed}=1
  \Longleftrightarrow
  y^{M}_{ed\mathrm{N}}=1,
  \ \rho^{M}_{e,d+1}=1,
  \ y^{M}_{e,d+2,\mathrm{A}}=1,
\]
並令
\[
  \Phi^{M,NA}=\sum_{e,d}\pi^{NA}_{ed}.
\]

\paragraph{Shift change without rest (\texttt{MS06})}
此規則判斷班別改變前是否出現實際休假。\(R\) 或 \(R1\) 會清除狀態，\(X\)
只被略過。若當日正常班別為 \(s\)，且
\(\gamma_{e,d-1}\notin\{0,s\}\)，則產生一次違反。總違反量記為
\[
  \Phi^{M,restswitch}.
\]
因此 \(\mathrm{E}\rightarrow X\rightarrow\mathrm{A}\) 仍屬未經休假換班；
\(\mathrm{E}\rightarrow R\rightarrow\mathrm{A}\) 則不違反。此規則與
\texttt{MS03} 不同：\texttt{MS03} 衡量同班別區塊長度，\texttt{MS06} 只判斷
換班前有沒有休假。

\paragraph{Preferred shift rotation (\texttt{MS07})}
令偏好轉移集合為
\[
  \mathcal{P}
  =
  \{(\mathrm{E},\mathrm{A}),(\mathrm{A},\mathrm{N}),(\mathrm{N},\mathrm{E})\}.
\]
依正常班別序列比較相鄰兩個有效正常班；\(R\)、\(R1\)、\(X\) 均不影響比較。
若班別改變且轉移不屬於 \(\mathcal{P}\)，則產生一次違反。總違反量記為
\[
  \Phi^{M,rotate}.
\]
此規則只判斷換班方向，與換班前是否休假無關。

\paragraph{Weekday rest fairness (\texttt{CS05})}
對每位人員定義目標月份平日休假數
\[
  R^{M,wd}_{e}
  =
  \sum_{d\in\mathcal{D}^{\mathrm{tar}}\cap\mathcal{D}^{wd}}
  \rho^{M}_{ed}.
\]
同一車站群組內的公平性違反量為
\[
  \Phi^{M,weekday}
  =
  \sum_{g\in\mathcal{G}}
  \left(
    \max_{e:g(h_e)=g}R^{M,wd}_{e}
    -
    \min_{e:g(h_e)=g}R^{M,wd}_{e}
  \right).
\]

\paragraph{Holiday rest fairness (\texttt{CS06})}
令
\[
  R^{M,hol}_{e}
  =
  \sum_{d\in\mathcal{D}^{\mathrm{tar}}\cap\mathcal{D}^{hol}}
  \rho^{M}_{ed}.
\]
則
\[
  \Phi^{M,holiday}
  =
  \sum_{g\in\mathcal{G}}
  \left(
    \max_{e:g(h_e)=g}R^{M,hol}_{e}
    -
    \min_{e:g(h_e)=g}R^{M,hol}_{e}
  \right).
\]

\paragraph{Cross-station support fairness (\texttt{MS08})}
令
\[
  S^{M}_{e}
  =
  \sum_{d\in\mathcal{D}^{\mathrm{tar}}}u_{ed}.
\]
則
\[
  \Phi^{M,support}
  =
  \sum_{g\in\mathcal{G}}
  \left(
    \max_{e:g(h_e)=g}S^{M}_{e}
    -
    \min_{e:g(h_e)=g}S^{M}_{e}
  \right).
\]

\subsubsection{Objective Groups}

M 模型定義五個目標群組：
\[
  J^M_1=\Phi^{M,R^*},
\]
\[
  J^M_2=\Phi^{M,ext},
\]
\[
  J^M_3
  =
  4\Phi^{M,month,R}
  +8\Phi^{M,month,R1},
\]
\[
  \begin{aligned}
    J^M_4={}&
    8\Phi^{M,home}
    +3\Phi^{M,streak}
    +2\Phi^{M,block}\\
    &+12\Phi^{M,NE}
    +8\Phi^{M,NA}
    +6\Phi^{M,restswitch},
  \end{aligned}
\]
\[
  J^M_5
  =
  \Phi^{M,rotate}
  +2\Phi^{M,weekday}
  +4\Phi^{M,holiday}
  +3\Phi^{M,support}.
\]
完整目標為
\[
  \operatorname{lexmin}
  \left(J^M_1,J^M_2,J^M_3,J^M_4,J^M_5\right).
\]

\subsection{The T Problem}

\subsubsection{Input and Output}

T 模型的輸入包括：
\begin{itemize}
  \item 目標月份與固定延伸天數 \(K\)；
  \item T 人員、專業組別、能力分數及建模日期的固定班組；
  \item 指定休假 \(R^*\)、固定勤務 \(X\) 與其他固定指派；
  \item 八週週期的一般 \(R\) 與 \(R1\) 額度；
  \item 目標月份的一般 \(R\) 與 \(R1\) 基準數量；
  \item 月初前連續七個日曆日的歷史班表與跨月夜轉早休假數；
  \item 建模開始前各人員、各八週週期已使用的 R 與 R1 數量；
  \item 國定假日集合。
\end{itemize}

模型輸出每位 T 人員在目標月份內的正常出勤、\(R\)、\(R1\) 或 \(X\)。正常
出勤班別必須等於該日期的固定班組。滿足 \(R^*\) 時，仍記錄底層實際使用的是
\(R\) 或 \(R1\)。延伸日期不列入正式輸出。

\subsubsection{Sets and Indices}

\begin{description}
  \item[\(\mathcal{E}^{T}\)] T 人員集合，索引為 \(e\)。
  \item[\(\mathcal{D}^{\mathrm{tar}}\)] 目標月份內的日期集合。
  \item[\(\mathcal{D}^{+}\)] 目標月份之後固定 \(K\) 日的延伸日期集合。
  \item[\(\mathcal{D}\)] 完整建模日期集合，
    \(\mathcal{D}=\mathcal{D}^{\mathrm{tar}}\cup\mathcal{D}^{+}\)。
  \item[\(\mathcal{S}\)] 班別集合，
    \(\mathcal{S}=\{\mathrm{E},\mathrm{A},\mathrm{N}\}\)。
  \item[\(\mathcal{C}\)] 八週休假週期集合。
  \item[\(\mathcal{Q}\)] 非空白專業組別集合，索引為 \(q\)。
  \item[\(\mathcal{E}_{ds}\)] 日期 \(d\) 固定屬於班別 \(s\) 的人員集合。
  \item[\(\mathcal{E}^{\mathrm{tar}}_s\)] 目標月份固定屬於班別 \(s\) 的人員集合。
  \item[\(\mathcal{S}^{\mathrm{tar}}\)] 目標月份實際有人員的班別集合。
  \item[\(\mathcal{E}^{N\rightarrow E}\)] 前月固定夜班、目標月固定早班的人員集合。
  \item[\(\mathcal{D}^{wd}\)] 星期一至星期五且不是國定假日的日期集合。
  \item[\(\mathcal{D}^{hol}\)] 星期六、星期日或國定假日的日期集合。
\end{description}

\subsubsection{Parameters}

令 \(\sigma_{ed}\in\mathcal{S}\) 表示人員 \(e\) 在日期 \(d\) 的固定班組。
目標月份與延伸日期都必須具有明確的 \(\sigma_{ed}\)。T 的班別時間為：早班
\(07{:}00\)--\(15{:}00\)、午班 \(15{:}00\)--\(23{:}00\)、夜班
\(23{:}00\)--翌日 \(07{:}00\)。夜班歸屬其開始日期。

令 \(\kappa_e\in\mathcal{Q}\cup\{\varnothing\}\) 為專業組別，令
\(\alpha_e\in\{1,2,3,4,5\}\) 為能力分數。依固定班組定義
\[
  \mathcal{E}_{ds}
  =
  \{e\in\mathcal{E}^{T}:\sigma_{ed}=s\},
  \qquad
  N_{ds}=|\mathcal{E}_{ds}|,
\]
以及每日正常出勤目標
\[
  B_{ds}=\left\lfloor\frac{N_{ds}}{2}\right\rfloor.
\]
班組內需要檢查的專業集合為
\[
  \mathcal{Q}_{ds}
  =
  \{q\in\mathcal{Q}:\exists e\in\mathcal{E}_{ds},\ \kappa_e=q\}.
\]

指定休假參數 \(p_{ed}\)、固定勤務參數 \(z_{ed}\)、八週額度
\(G_c,H_c\)、月休基準 \(n,m\) 與歷史額度
\(\overline{R}_{ec},\overline{R1}_{ec}\) 的意義與 M 模型相同。

令 \(d_0=\min\mathcal{D}^{\mathrm{tar}}-1\) 為目標月份前一日，令
\(d_1=\min\mathcal{D}^{\mathrm{tar}}\) 為目標月份第一日。對
\(e\in\mathcal{E}^{N\rightarrow E}\)，令 \(\beta_e\) 為最後一次實際夜班
結束後至 \(d_1\) 之前已出現的實際休假日數。月初前連續七日歷史另用來
重建法規連續日數與尚未結束的實際工作區段；公平性不使用歷史累積值。

\subsubsection{Variables}

主要決策變數為
\[
  x^{T}_{eds}\in\{0,1\},
\]
其中 \(x^{T}_{eds}=1\) 表示人員 \(e\) 在日期 \(d\) 正常出勤班別 \(s\)。
休假變數為
\[
  r^{T}_{ed},\ r^{T,1}_{ed}\in\{0,1\}.
\]
定義
\[
  v_{ed}=\sum_{s\in\mathcal{S}}x^{T}_{eds},
  \qquad
  w^{T}_{ed}=v_{ed}+z_{ed},
  \qquad
  \rho^{T}_{ed}=r^{T}_{ed}+r^{T,1}_{ed}.
\]
其中 \(X\) 是工作日，但不算正常出勤，因此不計入 T 的出勤、專業與能力指標。

令 \(q^{T,law}_{ed}\) 為只有一般 \(R\) 才歸零的法規連續日數；令
\(c^{T,work}_{ed}\) 為正常班或 \(X\) 形成的實際連續工作區段長度，\(R\) 與
\(R1\) 都會使其歸零。兩者用途不同，不得共用同一計數。

另定義：
\begin{itemize}
  \item \(\delta^{att}_{ds}\in\mathbb{Z}_{\geq0}\)：班組出勤不足量；
  \item \(\delta^{sp}_{dsq}\in\{0,1\}\)：某專業是否完全缺席；
  \item \(\delta^{ability}_{ds}\in\mathbb{Z}_{\geq0}\)：能力不足量；
  \item \(f_{ed}\in\{0,1\}\)：日期 \(d\) 是否為人員 \(e\) 在目標月份第一次正常早班；
  \item \(\delta^{bridge}_{ed}\in\mathbb{Z}_{\geq0}\)：夜班轉早班的休假不足量。
\end{itemize}

\subsubsection{Constraints}

\paragraph{Daily unique assignment (\texttt{CH01})}
\[
  \sum_{s\in\mathcal{S}}x^{T}_{eds}
  +r^{T}_{ed}+r^{T,1}_{ed}+z_{ed}=1,
  \qquad
  \forall e\in\mathcal{E}^{T},\ d\in\mathcal{D}.
\]
固定正常班、固定 \(R\)、固定 \(R1\) 與固定 \(X\) 直接固定對應變數。

\paragraph{Fixed monthly shift (\texttt{TH01})}
T 的正常出勤只能使用該日期固定班組。對所有 \(s\neq\sigma_{ed}\)，要求
\[
  x^{T}_{eds}=0.
\]
因此模型不能為了改善出勤、專業或能力而將人員跨班組調動。

\paragraph{Minimum eleven-hour rest (\texttt{CH03})}
令 \(\mathcal{I}^{T}\) 為所有實際休息時間少於十一小時的 T 正常班組合。對任一
\(((d,s),(d',s'))\in\mathcal{I}^{T}\)，要求
\[
  x^{T}_{eds}+x^{T}_{ed's'}\leq1,
  \qquad
  \forall e\in\mathcal{E}^{T}.
\]
正常班與固定 \(X\) 亦使用實際起訖時間檢查；兩筆固定 \(X\) 若彼此重疊或間隔
不足十一小時，輸入本身不可行。

\paragraph{At most six days without a general rest (\texttt{CH02})}
\[
  r^{T}_{ed}=1
  \ \Longrightarrow\ 
  q^{T,law}_{ed}=0,
\]
\[
  r^{T}_{ed}=0
  \ \Longrightarrow\ 
  q^{T,law}_{ed}=q^{T,law}_{e,d-1}+1,
\]
且
\[
  0\leq q^{T,law}_{ed}\leq6.
\]
正常班、\(X\) 與 \(R1\) 都會使計數增加；只有一般 \(R\) 會歸零。

\paragraph{Eight-week rest quotas (\texttt{CH04})}
對每個週期 \(c\)，令
\[
  \mathcal{D}_c=c\cap\mathcal{D},
  \qquad
  F_c=\left|\{d\in c:d>\max\mathcal{D}\}\right|.
\]
若模型已涵蓋週期結束日，則
\[
  \overline{R}_{ec}+\sum_{d\in\mathcal{D}_c}r^{T}_{ed}=G_c,
  \qquad
  \overline{R1}_{ec}+\sum_{d\in\mathcal{D}_c}r^{T,1}_{ed}=H_c.
\]
若尚未涵蓋週期結束日，則
\[
  \overline{R}_{ec}+\sum_{d\in\mathcal{D}_c}r^{T}_{ed}\leq G_c,
  \qquad
  \overline{R1}_{ec}+\sum_{d\in\mathcal{D}_c}r^{T,1}_{ed}\leq H_c,
\]
\[
  \overline{R}_{ec}+\sum_{d\in\mathcal{D}_c}r^{T}_{ed}+F_c\geq G_c,
\]
\[
  \overline{R1}_{ec}+\sum_{d\in\mathcal{D}_c}r^{T,1}_{ed}+F_c\geq H_c,
\]
以及
\[
  \overline{R}_{ec}+\overline{R1}_{ec}
  +\sum_{d\in\mathcal{D}_c}\left(r^{T}_{ed}+r^{T,1}_{ed}\right)
  +F_c
  \geq G_c+H_c.
\]

\paragraph{Fixed assignments and historical boundary (\texttt{CH05})}
所有固定指派均不得由模型改寫。歷史班表決定建模區間第一日前的
\(q^{T,law}\)、\(c^{T,work}\)、八週累積值與跨月夜轉早所需狀態。

\paragraph{Requested rest (\texttt{CS01})}
\[
  \Phi^{T,R^*}
  =
  \sum_{\substack{e\in\mathcal{E}^{T},\ d\in\mathcal{D}^{\mathrm{tar}}\\p_{ed}=1}}
  \left(1-r^{T}_{ed}-r^{T,1}_{ed}\right).
\]

\paragraph{Monthly rest distribution (\texttt{CS02}, \texttt{CS03})}
令
\[
  Q^{T,R}_e=\sum_{d\in\mathcal{D}^{\mathrm{tar}}}r^{T}_{ed},
  \qquad
  Q^{T,R1}_e=\sum_{d\in\mathcal{D}^{\mathrm{tar}}}r^{T,1}_{ed}.
\]
沿用
\[
  f_t(q)=\left[\max\{0,|q-t|-1\}\right]^2,
\]
則
\[
  \Phi^{T,month,R}
  =
  \sum_{e\in\mathcal{E}^{T}}f_n\!\left(Q^{T,R}_e\right),
\]
\[
  \Phi^{T,month,R1}
  =
  \sum_{e\in\mathcal{E}^{T}}f_m\!\left(Q^{T,R1}_e\right).
\]

\paragraph{Shift attendance (\texttt{TS01})}
日期 \(d\)、班別 \(s\) 的正常出勤人數為
\[
  A_{ds}
  =
  \sum_{e\in\mathcal{E}_{ds}}x^{T}_{eds}.
\]
要求
\[
  \delta^{att}_{ds}\geq B_{ds}-A_{ds},
  \qquad
  \delta^{att}_{ds}\geq0.
\]
總違反量為
\[
  \Phi^{T,att}
  =
  \sum_{d\in\mathcal{D}^{\mathrm{tar}}}
  \sum_{s\in\mathcal{S}}\delta^{att}_{ds}.
\]

\paragraph{Specialty attendance (\texttt{TS02})}
對每個 \(q\in\mathcal{Q}_{ds}\)，令
\[
  A^{sp}_{dsq}
  =
  \sum_{\substack{e\in\mathcal{E}_{ds}\\\kappa_e=q}}x^{T}_{eds},
\]
並令
\[
  M^{sp}_{dsq}
  =
  \left|\{e\in\mathcal{E}_{ds}:\kappa_e=q\}\right|.
\]
以 \(\delta^{sp}_{dsq}=1\) 表示該專業無人正常出勤：
\[
  A^{sp}_{dsq}\geq1-\delta^{sp}_{dsq},
  \qquad
  A^{sp}_{dsq}\leq M^{sp}_{dsq}(1-\delta^{sp}_{dsq}).
\]
總違反量為
\[
  \Phi^{T,specialty}
  =
  \sum_{d\in\mathcal{D}^{\mathrm{tar}}}
  \sum_{s\in\mathcal{S}}
  \sum_{q\in\mathcal{Q}_{ds}}\delta^{sp}_{dsq}.
\]
專業空白的人員仍可正常出勤，但空白不形成需要覆蓋的專業類別。

\paragraph{Ability level (\texttt{TS03})}
日期 \(d\)、班別 \(s\) 的能力總和為
\[
  C_{ds}
  =
  \sum_{e\in\mathcal{E}_{ds}}\alpha_e x^{T}_{eds}.
\]
正常出勤者的平均能力希望至少為 3，因此
\[
  \delta^{ability}_{ds}
  \geq
  3A_{ds}-C_{ds},
  \qquad
  \delta^{ability}_{ds}\geq0.
\]
總違反量為
\[
  \Phi^{T,ability}
  =
  \sum_{d\in\mathcal{D}^{\mathrm{tar}}}
  \sum_{s\in\mathcal{S}}\delta^{ability}_{ds}.
\]
若無人正常出勤，能力不足量為 0；無人出勤由 \texttt{TS01} 處理。

\paragraph{Continuous work-streak length (\texttt{CS04})}
T 的實際工作區段同樣由連續正常班或 \(X\) 組成，\(R\) 與 \(R1\) 都會結束
區段。沿用 M 模型的 \(D(L)\)。令 \(\mathcal{B}^{T,work}_e\) 為承接歷史後的
最大實際工作區段集合，則
\[
  \Phi^{T,streak}
  =
  \sum_{e\in\mathcal{E}^{T}}
  \sum_{B\in\mathcal{B}^{T,work}_e:\ \max B\in\mathcal{D}^{\mathrm{tar}}}
  D(|B|).
\]
此規則控制合理工作區段長度，與只有一般 \(R\) 才歸零的 \texttt{CH02} 不同。

\paragraph{Night-to-early month transition rest (\texttt{TS04})}
此規則只適用於 \(e\in\mathcal{E}^{N\rightarrow E}\)。令 \(f_{ed}=1\) 表示日期
\(d\) 是該人在目標月份第一次正常早班。若 \(f_{ed}=1\)，最後夜班至該日之前
的休假日數為
\[
  B^{bridge}_{ed}
  =
  \beta_e
  +
  \sum_{\substack{\tau\in\mathcal{D}^{\mathrm{tar}}\\\tau<d}}
  \rho^{T}_{e\tau}.
\]
定義
\[
  \delta^{bridge}_{ed}
  =
  \begin{cases}
    \max(0,2-B^{bridge}_{ed}), & f_{ed}=1,\\
    0, & f_{ed}=0.
  \end{cases}
\]
總違反量為
\[
  \Phi^{T,monthrest}
  =
  \sum_{e\in\mathcal{E}^{N\rightarrow E}}
  \sum_{d\in\mathcal{D}^{\mathrm{tar}}}\delta^{bridge}_{ed}.
\]
此規則要求夜班轉早班的人員，在最後夜班與第一次早班之間至少有兩個實際休假日。

\paragraph{Month-boundary rest balance (\texttt{TS05})}
令 \(\overline{\rho}^{T}_{e,d_0}\) 為前月最後一日的歷史休假指示，並定義
\[
  B^{-}
  =
  \sum_{e\in\mathcal{E}^{N\rightarrow E}}
  \overline{\rho}^{T}_{e,d_0},
  \qquad
  B^{+}
  =
  \sum_{e\in\mathcal{E}^{N\rightarrow E}}
  \rho^{T}_{e,d_1}.
\]
月交界休假分散違反量為
\[
  \Phi^{T,monthbalance}=|B^{-}-B^{+}|.
\]

\paragraph{Weekday rest fairness (\texttt{CS05})}
令
\[
  R^{T,wd}_{e}
  =
  \sum_{d\in\mathcal{D}^{\mathrm{tar}}\cap\mathcal{D}^{wd}}
  \rho^{T}_{ed}.
\]
只在目標月份同一固定班組內比較：
\[
  \Phi^{T,weekday}
  =
  \sum_{s\in\mathcal{S}^{\mathrm{tar}}}
  \left(
    \max_{e\in\mathcal{E}^{\mathrm{tar}}_s}R^{T,wd}_{e}
    -
    \min_{e\in\mathcal{E}^{\mathrm{tar}}_s}R^{T,wd}_{e}
  \right).
\]

\paragraph{Holiday rest fairness (\texttt{CS06})}
令
\[
  R^{T,hol}_{e}
  =
  \sum_{d\in\mathcal{D}^{\mathrm{tar}}\cap\mathcal{D}^{hol}}
  \rho^{T}_{ed}.
\]
則
\[
  \Phi^{T,holiday}
  =
  \sum_{s\in\mathcal{S}^{\mathrm{tar}}}
  \left(
    \max_{e\in\mathcal{E}^{\mathrm{tar}}_s}R^{T,hol}_{e}
    -
    \min_{e\in\mathcal{E}^{\mathrm{tar}}_s}R^{T,hol}_{e}
  \right).
\]

\subsubsection{Objective Groups}

T 模型定義五個目標群組：
\[
  J^T_1=\Phi^{T,R^*},
\]
\[
  J^T_2
  =
  9\Phi^{T,att}
  +3\Phi^{T,specialty}
  +\Phi^{T,ability},
\]
\[
  J^T_3
  =
  \Phi^{T,month,R}
  +\Phi^{T,month,R1},
\]
\[
  J^T_4
  =
  3\Phi^{T,streak}
  +12\Phi^{T,monthrest}
  +5\Phi^{T,monthbalance},
\]
\[
  J^T_5
  =
  2\Phi^{T,weekday}
  +4\Phi^{T,holiday}.
\]
完整目標為
\[
  \operatorname{lexmin}
  \left(J^T_1,J^T_2,J^T_3,J^T_4,J^T_5\right).
\]

\section{Alternative Candidates}

只有在五個目標群組皆證明為最優後，系統才搜尋替代候選。每個模型最多輸出
三份候選；所有已固定的 \(J_1,\ldots,J_5\) 最優值維持不變。

候選差異只比較目標月份中真正由求解器決定的 \((e,d)\) 格。固定正常班、
固定 \(R\)、固定 \(R1\)、固定 \(X\) 與延伸日期均排除；指定休假 \(R^*\)
仍由求解器決定底層使用 \(R\) 或 \(R1\)，因此保留在比較集合。比較實際狀態
時忽略 \(R^*\) 顯示標記，但 \(R\) 與 \(R1\) 仍視為不同狀態。

若可比較格數為 \(N_{cell}\)，定義最低差異格數
\[
  \Delta_{min}=\left\lceil0.1N_{cell}\right\rceil.
\]
第二份候選與第一份至少相差 \(\Delta_{min}\) 格；第三份必須分別與前兩份
達到相同門檻。候選搜尋共用整次求解的剩餘時間；時間用完時保留已找到的候選。

\section{Mathematical Symbols and C\# Variables}

本章是數學模型與 source-as-spec 實作的直接對照。Rule ID 不寫入程式；以下
函數名稱就是公式在原始碼中的定位方式。公開資料型別位於
\texttt{SolverContracts.cs}、\texttt{MContracts.cs} 與
\texttt{TContracts.cs}，建模函數分別只位於 \texttt{MSolver.cs} 與
\texttt{TSolver.cs}。

\subsection{Input and Sets}

\small
\begin{longtable}{p{2.7cm}p{5.2cm}p{6.4cm}}
  \toprule
  數學符號 & C\# 資料 & 來源或建立位置 \\
  \midrule
  \endfirsthead
  \toprule
  數學符號 & C\# 資料 & 來源或建立位置 \\
  \midrule
  \endhead
  \(\mathcal D^{\mathrm{tar}}\) & \texttt{DateOnly[] targetDates} & \texttt{TargetDates(input)} \\
  \(\mathcal D^{+},\mathcal D\) & \texttt{DateOnly[] modelDates} & \texttt{ModelDates(input)}；目標月後 \texttt{ExtensionDays} 日 \\
  \(\mathcal E^M\) & \texttt{MInput.Employees} & \texttt{IReadOnlyList<MEmployee>} \\
  \(\mathcal E^T\) & \texttt{TInput.Employees} & \texttt{IReadOnlyList<TEmployee>} \\
  \(\mathcal C,G_c,H_c\) & \texttt{RestCycles} & \texttt{RestCycle.Id/Start/End/}\newline\texttt{RequiredRest/RequiredSpecialRest} \\
  \(\overline R_{ec},\overline{R1}_{ec}\) & \texttt{UsageBeforeMonth} & \texttt{EmployeeCycleUsage.}\newline\texttt{RestUsed/SpecialRestUsed} \\
  \(n,m\) & \texttt{MonthlyRestTarget}, \texttt{MonthlySpecialRestTarget} & \texttt{MInput} 或 \texttt{TInput} \\
  \(h_e\) & \texttt{MEmployee.HomeStation} & \texttt{MContracts.cs} \\
  \(p_{ed}\) & \texttt{RequestedRests} & \texttt{RequestedRest(EmployeeId, Date)} \\
  \(z_{ed}\) & \texttt{WorkEvents} & 有相同人員與日期的 \texttt{WorkEvent} 即為 1 \\
  \(\sigma_{ed}\) & \texttt{TInput.FixedShifts} & \texttt{TFixedShift(EmployeeId, Date, Shift)} \\
  \(\kappa_e,\alpha_e\) & \texttt{TEmployee.Specialty/Ability} & \texttt{TContracts.cs} \\
  \(\beta_e\) & \texttt{NightToEarlyHistory.}\newline\texttt{RestDaysAfterLastNight} & 只供前月夜班轉目標月早班者 \\
  月初七日歷史 & \texttt{PreviousSevenDays} & \texttt{MHistoryDay} 或 \texttt{THistoryDay}；重建序列狀態，不重算八週額度 \\
  \bottomrule
\end{longtable}
\normalsize

\subsection{M Model}

\small
\begin{longtable}{p{2.7cm}p{5.2cm}p{6.4cm}}
  \toprule
  數學符號 & C\# 變數 & 建立或使用函數 \\
  \midrule
  \endfirsthead
  \toprule
  數學符號 & C\# 變數 & 建立或使用函數 \\
  \midrule
  \endhead
  \(x^M_{edls}\) & \texttt{variables.Work[(e,d,l,s)]}，\texttt{BoolVar} & \texttt{CreateVariables} \\
  \(r^M_{ed}\) & \texttt{variables.Rest[(e,d)]}，\texttt{BoolVar} & \texttt{CreateVariables} \\
  \(r^{M,1}_{ed}\) & \texttt{variables.SpecialRest}\newline\texttt{[(e,d)]}，\texttt{BoolVar} & \texttt{CreateVariables} \\
  \(a_{dls}\) & \texttt{variables.External[(d,l,s)]}，\texttt{IntVar} & \texttt{CreateVariables}、\texttt{AddExactCoverage} \\
  \(y^M_{eds}\) & \texttt{variables.WorksShift}\newline\texttt{[(e,d,s)]}，\texttt{BoolVar} & \texttt{CreateVariables} \\
  \(w^M_{ed}\) & \texttt{variables.ActualWork}\newline\texttt{[(e,d)]}，\texttt{BoolVar} & 正常班或 \(X\) 為 1 \\
  \(\rho^M_{ed}\) & \texttt{variables.AnyRest[(e,d)]}，\texttt{BoolVar} & \(R+R1\) \\
  \(u_{ed}\) & \texttt{variables.}\newline\texttt{SupportsOtherStation[(e,d)]}，\texttt{BoolVar} & 非本站正常班之和 \\
  \(q^{law}_{ed}\) & 七日滑動視窗，未另建狀態 class & \texttt{AddMaximumSixDays}\newline\texttt{WithoutGeneralRest} \\
  \(c^{work}_{ed}\) & 函數內 \texttt{IntVar count} & \texttt{WorkStreakViolation} \\
  \((\lambda_{ed},\ell_{ed})\) & 函數內 \texttt{lastShift/blockLength} & \texttt{SameShiftBlockViolation} \\
  \(\gamma_{ed}\) & 函數內 \texttt{state} & \texttt{ShiftChangeWithoutRestViolation} \\
  \(\eta_{ed}\) & 函數內 \texttt{state} & \texttt{NonPreferredRotationViolation} \\
  \(J^M_1,\ldots,J^M_5\) & \texttt{objectives} 中的 \texttt{LinearExpr Total} & \texttt{BuildObjectives}；權重直接寫在同一函數 \\
  \bottomrule
\end{longtable}
\normalsize

\subsection{T Model and Solver Control}

\footnotesize
\begin{longtable}{p{2.7cm}p{5.2cm}p{6.4cm}}
  \toprule
  數學符號或概念 & C\# 變數 & 建立或使用函數 \\
  \midrule
  \endfirsthead
  \toprule
  數學符號或概念 & C\# 變數 & 建立或使用函數 \\
  \midrule
  \endhead
  \(x^T_{eds}\) & \texttt{variables.Work[(e,d,s)]}，\texttt{BoolVar} & \texttt{CreateVariables} \\
  \(r^T_{ed},r^{T,1}_{ed}\) & \texttt{variables.Rest/SpecialRest}\newline\texttt{[(e,d)]} & \texttt{CreateVariables} \\
  \(w^T_{ed},\rho^T_{ed}\) & \texttt{variables.ActualWork/}\newline\texttt{AnyRest[(e,d)]} & \texttt{CreateVariables} \\
  \(\delta^{att}_{ds}\) & 函數內 \texttt{IntVar deficit} & \texttt{AttendanceViolation} \\
  \(\delta^{sp}_{dsq}\) & 函數內 \texttt{BoolVar absent} & \texttt{SpecialtyViolation} \\
  \(\delta^{ability}_{ds}\) & 函數內 \texttt{IntVar deficit} & \texttt{AbilityViolation} \\
  \(f_{ed},\delta^{bridge}_{ed}\) & 函數內 \texttt{first/deficit} & \texttt{NightToEarlyRestViolation} \\
  \(J^T_1,\ldots,J^T_5\) & \texttt{objectives} 中的 \texttt{LinearExpr Total} & \texttt{BuildObjectives}；權重直接寫在同一函數 \\
  字典序固定 & \texttt{model.Add(objective.Total == optimum)} & \texttt{TSolver.Solve} 或 \texttt{MSolver.Solve} \\
  總求解時限 & \texttt{SolverOptions.TimeLimit} & 每次 \texttt{Solve} 只使用剩餘秒數 \\
  主動取消 & \texttt{CancellationToken} & 註冊 \texttt{CpSolver.StopSearch}；取消後拋出 \texttt{OperationCanceledException} \\
  \bottomrule
\end{longtable}
\normalsize

% -----------------------------------------------------------
%    You don't have to mess with anything below this line.
% -----------------------------------------------------------
\end{document}
